% =====================================================================
%  CARNET D'ENTRAÎNEMENT — FICHE 6 : Les limites
%  Thème 1 — CORRIGÉ — Niveau DIFFICILE
% =====================================================================
\documentclass[11pt,a4paper]{article}
\usepackage[utf8]{inputenc}
\usepackage[T1]{fontenc}
\usepackage{lmodern}
\usepackage[french]{babel}
\usepackage{amsmath,amssymb,mathtools}
\usepackage{icomma}
\usepackage{xcolor}
\usepackage{tikz}
\usepackage{pgfplots}
\usepackage[most]{tcolorbox}
\usepackage{enumitem}
\usepackage{array}
\usepackage{booktabs}
\usepackage{fancyhdr}
\usepackage[hidelinks]{hyperref}
\usetikzlibrary{arrows.meta,calc,intersections}
\pgfplotsset{compat=1.18}
\usepackage[a4paper,margin=2.2cm,headheight=26pt,headsep=10pt]{geometry}

\definecolor{primary}{RGB}{223,87,99}
\definecolor{primarydark}{RGB}{150,42,55}
\definecolor{excol}{RGB}{0,135,90}
\definecolor{curvecol}{RGB}{40,80,190}

\tcbset{boxbase/.style={enhanced,breakable,boxrule=0.9pt,arc=2.5pt,
   left=3mm,right=3mm,top=2.3mm,bottom=2.3mm,
   fonttitle=\bfseries,coltitle=white,toptitle=0.6mm,bottomtitle=0.6mm}}
\newtcolorbox[auto counter]{correction}[1][]{boxbase,colback=excol!5,colframe=excol,
  title={\textsc{Corrigé}~\thetcbcounter\ifx\relax#1\relax\relax\else\textmd{ --- \itshape #1}\fi}}

\pagestyle{fancy}
\fancyhf{}
\fancyhead[L]{\small\textcolor{primarydark}{\textbf{Fiche 6} — Les limites}}
\fancyhead[R]{\small\textcolor{primarydark}{\textsc{Thème 1 — Corrigé Difficile}}}
\fancyfoot[C]{\small\thepage}
\renewcommand{\headrulewidth}{0.8pt}
\renewcommand{\headrule}{\hbox to\headwidth{\color{primary}\leaders\hrule height \headrulewidth\hfill}}

\newcommand{\R}{\mathbb{R}}

\begin{document}

\begin{center}
{\Large\bfseries\color{primarydark}Thème 1 — Notations, lecture \& calcul direct}\\[2pt]
{\large\color{excol}Corrigé — Niveau Difficile}
\end{center}
\vspace{6pt}

\begin{correction}[une fonction à deux paramètres]
\begin{enumerate}[label=\alph*),leftmargin=2em,itemsep=3pt]
  \item À gauche ($x<2$) : $f(x)=ax+1$, donc $\displaystyle\lim_{x\to 2^{-}} f(x)=2a+1$.\\
        À droite ($x>2$) : $f(x)=x^2+b$, donc $\displaystyle\lim_{x\to 2^{+}} f(x)=4+b$.
  \item $f$ \emph{possède une limite} en $2$ ssi les deux limites latérales sont égales :
        \[ 2a+1=4+b\quad\Longleftrightarrow\quad \boxed{b=2a-3}. \]
  \item Pour la \emph{continuité} en $2$, il faut de plus que cette limite commune soit égale à $f(2)=3$ :
        \[ 2a+1=3\ \text{(et)}\ 4+b=3. \]
  \item De $2a+1=3$ on tire $a=1$ ; de $4+b=3$ on tire $b=-1$. (Ces valeurs vérifient bien $b=2a-3=-1$.)
        \[ \boxed{(a,b)=(1,-1)} \]
        est l'unique paire rendant $f$ continue en $2$. Exemple : $f(x)=x+1$ pour $x<2$, $f(2)=3$,
        $f(x)=x^2-1$ pour $x>2$ ; on a alors $\lim_{x\to2^\pm} f=3=f(2)$.
\end{enumerate}
\end{correction}

\begin{correction}[lecture graphique complète]
\begin{enumerate}[label=\alph*),leftmargin=2em,itemsep=3pt]
  \item Aux deux bouts la courbe colle à la droite $y=1$ :
        $\displaystyle\lim_{x\to +\infty} f(x)=\lim_{x\to -\infty} f(x)=1$. La droite $y=1$ est une
        \textbf{asymptote horizontale}.
  \item À gauche de $2$, la branche plonge : $\displaystyle\lim_{x\to 2^{-}} f(x)=-\infty$. La droite
        $x=2$ est donc une \textbf{asymptote verticale}.
  \item Le rond ouvert est en $(2\,;3)$, donc $\displaystyle\lim_{x\to 2^{+}} f(x)=3$ ; le rond plein est
        en $(2\,;0)$, donc $f(2)=0$.
  \item Les limites latérales en $2$ diffèrent ($-\infty$ à gauche, $3$ à droite) : $f$ \textbf{n'a pas de
        limite} en $2$ et n'est donc pas continue. De plus, à droite $\lim_{x\to 2^{+}} f(x)=3\neq f(2)=0$ :
        $f$ \textbf{n'est pas continue à droite} non plus. (La valeur $f(2)=0$ est un point isolé.)
\end{enumerate}
\end{correction}

\begin{correction}[mise en situation : un tarif par paliers]
\begin{enumerate}[label=\alph*),leftmargin=2em,itemsep=3pt]
  \item À gauche ($t\leq 20$) : $M(t)=30$, donc $\displaystyle\lim_{t\to 20^{-}} M(t)=30$.\\
        À droite ($t>20$) : $M(t)=30+2(t-20)$, donc $\displaystyle\lim_{t\to 20^{+}} M(t)=30+2(20-20)=30$.
  \item Les deux limites valent $30$ et $M(20)=30$ : \textbf{$M$ est continue en $20$}. Concrètement, le tarif
        \emph{ne saute pas} : à $20$ minutes on paie $30$\,€, et juste après le montant augmente
        continûment (pas de bond brutal).
  \item Avec $35+2(t-20)$ pour $t>20$ : $\displaystyle\lim_{t\to 20^{+}} M(t)=35$ tandis que
        $\displaystyle\lim_{t\to 20^{-}} M(t)=30$. Les limites latérales diffèrent : $M$ serait
        \textbf{discontinue} en $20$, avec un \textbf{saut de $35-30=5$\,€} (le prix bondirait de $5$\,€
        dès qu'on dépasse $20$ minutes).
\end{enumerate}
\end{correction}

\end{document}
